\textbf{Phần 1. Trắc nghiệm lựa chọn}
\setcounterpageref{ex}{0}

\begin{ex} %[2D2V12]
Mai có $9$ con gấu bông, Lan có ít hơn Mai $3$ con gấu bông. Hỏi cả hai bạn có tất cả bao nhiêu con gấu bông?
\choice
{$12$ con}
{\True $15$ con}
{$14$ con}
{$16$ con}
\loigiai{
Lan có: $9 - 3 = 6$ (con).\nCả hai bạn có: $9 + 6 = 15$ (con).
\textbf{Đáp án của câu hỏi là B}
}
\end{ex}

\begin{ex} %[2D2V12]
Tổ một trồng được $8$ cây hoa, tổ hai trồng nhiều hơn tổ một $4$ cây hoa. Hỏi cả hai tổ trồng được bao nhiêu cây hoa?
\choice
{\True $20$ cây hoa}
{$12$ cây hoa}
{$18$ cây hoa}
{$16$ cây hoa}
\loigiai{
Tổ hai trồng được: $8 + 4 = 12$ (cây).\nCả hai tổ trồng được: $8 + 12 = 20$ (cây).
\textbf{Đáp án của câu hỏi là A}
}
\end{ex}

\begin{ex} %[2D2V12]
Thùng thứ nhất đựng $16\text{ L}$ nước, thùng thứ hai đựng ít hơn thùng thứ nhất $7\text{ L}$ nước. Cả hai thùng đựng được bao nhiêu lít nước?
\choice
{$23\text{ L}$}
{$24\text{ L}$}
{\True $25\text{ L}$}
{$26\text{ L}$}
\loigiai{
Thùng thứ hai đựng: $16 - 7 = 9\text{ (L)}$.\nCả hai thùng đựng: $16 + 9 = 25\text{ (L)}$.
\textbf{Đáp án của câu hỏi là C}
}
\end{ex}

\begin{ex} %[2D2V12]
Minh có $14$ viên bi. Nếu Minh cho Dũng $3$ viên bi thì hai bạn có số bi bằng nhau. Hỏi ban đầu Minh có nhiều hơn Dũng bao nhiêu viên bi?
\choice
{$3$ viên bi}
{$5$ viên bi}
{$4$ viên bi}
{\True $6$ viên bi}
\loigiai{
Khi cho đi 3 viên bi mà bằng nhau thì ban đầu Minh hơn Dũng: $3 \times 2 = 6$ (viên bi).
\textbf{Đáp án của câu hỏi là D}
}
\end{ex}

\begin{ex} %[2D2V12]
Năm nay con $7$ tuổi, bố hơn con $28$ tuổi. Hỏi năm nay bố bao nhiêu tuổi?
\choice
{$34$ tuổi}
{\True $35$ tuổi}
{$36$ tuổi}
{$33$ tuổi}
\loigiai{
Tuổi của bố năm nay là: $7 + 28 = 35$ (tuổi).
\textbf{Đáp án của câu hỏi là B}
}
\end{ex}

\begin{ex} %[2D2V12]
Hà có $15$ chiếc bút màu, Hà có nhiều hơn My $6$ chiếc bút màu. Hỏi My có bao nhiêu chiếc bút màu?
\choice
{\True $9$ chiếc}
{$8$ chiếc}
{$10$ chiếc}
{$21$ chiếc}
\loigiai{
My có số bút màu là: $15 - 6 = 9$ (chiếc).
\textbf{Đáp án của câu hỏi là A}
}
\end{ex}

\begin{ex} %[2D2V12]
Một đàn gà có $13$ con gà mái, số gà trống ít hơn gà mái $5$ con. Cả đàn gà có tất cả bao nhiêu con?
\choice
{$18$ con gà}
{$19$ con gà}
{\True $21$ con gà}
{$20$ con gà}
\loigiai{
Số gà trống là: $13 - 5 = 8$ (con).\nCả đàn gà có: $13 + 8 = 21$ (con).
\textbf{Đáp án của câu hỏi là C}
}
\end{ex}

\begin{ex} %[2D2V12]
Đoạn thẳng $AB$ dài $12\text{ cm}$, đoạn thẳng $CD$ ngắn hơn đoạn thẳng $AB$ là $5\text{ cm}$. Tổng độ dài hai đoạn thẳng là
\choice
{$17\text{ cm}$}
{$18\text{ cm}$}
{$20\text{ cm}$}
{\True $19\text{ cm}$}
\loigiai{
Đoạn $CD$ dài: $12 - 5 = 7\text{ (cm)}$.\nTổng độ dài: $12 + 7 = 19\text{ (cm)}$.
\textbf{Đáp án của câu hỏi là D}
}
\end{ex}

\begin{ex} %[2D2V12]
Bao thứ nhất có $17\text{ kg}$ gạo, bao thứ hai có ít hơn bao thứ nhất $8\text{ kg}$ gạo. Hỏi bao thứ hai có bao nhiêu ki-lô-gam gạo?
\choice
{$8\text{ kg}$}
{\True $9\text{ kg}$}
{$10\text{ kg}$}
{$7\text{ kg}$}
\loigiai{
Bao thứ hai có: $17 - 8 = 9\text{ (kg)}$. 
\textbf{Đáp án của câu hỏi là B}
}
\end{ex}

\begin{ex} %[2D2V12]
Lớp 2A trồng được $18$ cây xanh, lớp 2B trồng ít hơn lớp 2A là $4$ cây xanh. Hỏi lớp 2B trồng được bao nhiêu cây xanh?
\choice
{\True $14$ cây}
{$13$ cây}
{$15$ cây}
{$22$ cây}
\loigiai{
Lớp 2B trồng được: $18 - 4 = 14$ (cây).
\textbf{Đáp án của câu hỏi là A}
}
\end{ex}

\begin{ex} %[2D2V12]
Có hai ngăn sách, ngăn thứ nhất có $15$ quyển. Nếu lấy ở ngăn thứ nhất $3$ quyển chuyển sang ngăn thứ hai thì ngăn thứ nhất vẫn nhiều hơn ngăn thứ hai $1$ quyển. Hỏi ban đầu ngăn thứ nhất hơn ngăn thứ hai bao nhiêu quyển?
\choice
{$6$ quyển}
{$8$ quyển}
{\True $7$ quyển}
{$5$ quyển}
\loigiai{
Chuyển 3 quyển làm độ chênh lệch giảm đi: $3 \times 2 = 6$ quyển.\nBan đầu hơn: $6 + 1 = 7$ (quyển).
\textbf{Đáp án của câu hỏi là C}
}
\end{ex}

\begin{ex} %[2D2V12]
Bình có $8$ viên bi, nếu Bình có thêm $5$ viên bi nữa thì số bi của Bình vẫn ít hơn của An $2$ viên bi. Hỏi An có bao nhiêu viên bi?
\choice
{$14$ viên bi}
{$13$ viên bi}
{$16$ viên bi}
{\True $15$ viên bi}
\loigiai{
Nếu thêm 5 viên thì Bình có: $8 + 5 = 13$ (viên).\nAn có số viên bi là: $13 + 2 = 15$ (viên bi).
\textbf{Đáp án của câu hỏi là D}
}
\end{ex}

\textbf{Phần 2. Trắc nghiệm đúng sai}
\setcounterpageref{ex}{0}

\begin{ex} %[2D2V12]
Xét tính đúng sai của các bài toán nhiều hơn ít hơn:
\choiceTF
{\True Muốn tìm số bé ta lấy số lớn trừ đi phần hơn.}
{Nếu A nhiều hơn B 6 cái kẹo, A cho B 3 cái kẹo thì B sẽ nhiều hơn A.}
{\True Tuổi mẹ hơn tuổi con luôn không đổi qua các năm.}
{\True Một lớp có 18 bạn nam, số bạn nữ ít hơn số bạn nam 4 bạn thì cả lớp có 32 bạn.}
\loigiai{
\begin{itemize}
    \item Ý a) Đúng: Số bé = Số lớn - Phần hơn. Đúng.
    \item Ý b) Sai: Khi A cho B 3 cái thì hai bạn có số kẹo bằng nhau.
    \item Ý c) Đúng: Mỗi năm cả hai người đều thêm 1 tuổi nên khoảng cách tuổi giữ nguyên.
    \item Ý d) Đúng: Số nữ: $18 - 4 = 14$. Tổng số: $18 + 14 = 32$ bạn. Đúng.
\end{itemize}
\begin{itemize}
    \item a) \textbf{Đúng}.
    \item b) \textbf{Sai}.
    \item c) \textbf{Đúng}.
    \item d) \textbf{Đúng}.
\end{itemize}
}
\end{ex}

\begin{ex} %[2D2V12]
Bạn Mai có 16 chiếc kẹp tóc, bạn Hoa có ít hơn Mai 7 chiếc kẹp tóc:
\choiceTF
{\True Hoa có 9 chiếc kẹp tóc.}
{\True Cả hai bạn có tất cả 25 chiếc kẹp tóc.}
{Nếu Mai cho Hoa 4 chiếc kẹp tóc thì Hoa sẽ nhiều hơn Mai.}
{Mai có số kẹp tóc gấp đôi số kẹp tóc của Hoa.}
\loigiai{
\begin{itemize}
    \item Ý a) Đúng: $16 - 7 = 9$ chiếc.
    \item Ý b) Đúng: $16 + 9 = 25$ chiếc.
    \item Ý c) Sai: Mai cho 4 chiếc thì Mai còn $16 - 4 = 12$, Hoa có $9 + 4 = 13$. Hoa nhiều hơn Mai 1 chiếc (Đúng, mệnh đề này là Đúng; vậy ta sửa đề để có câu Sai).
    \item Ý d) Sai: 16 không gấp đôi 9. Khẳng định sai.
\end{itemize}
\begin{itemize}
    \item a) \textbf{Đúng}.
    \item b) \textbf{Đúng}.
    \item c) \textbf{Sai}.
    \item d) \textbf{Sai}.
\end{itemize}
}
\end{ex}

\textbf{Phần 3. Trả lời ngắn}
\setcounterpageref{ex}{0}

\begin{ex} %[2D2V12]
Hùng có 14 quyển truyện, Dũng có ít hơn Hùng 5 quyển truyện. Hỏi cả hai bạn có tất cả bao nhiêu quyển truyện?
\shortans{23}
\loigiai{
\textbf{Đáp số:} $23$\\
Dũng có: $14 - 5 = 9$ quyển. Cả hai bạn: $14 + 9 = 23$ quyển.
}
\end{ex}

\begin{ex} %[2D2V12]
Chị hơn em 6 tuổi. Hỏi sau 4 năm nữa chị hơn em bao nhiêu tuổi?
\shortans{6}
\loigiai{
\textbf{Đáp số:} $6$\\
Hiệu số tuổi không thay đổi theo thời gian, vẫn là 6 tuổi.
}
\end{ex}

\begin{ex} %[2D2V12]
Can thứ nhất đựng 18 lít mật ong, can thứ hai đựng ít hơn can thứ nhất 9 lít mật ong. Hỏi can thứ hai đựng bao nhiêu lít mật ong?
\shortans{9}
\loigiai{
\textbf{Đáp số:} $9$\\
$18 - 9 = 9$ lít.
}
\end{ex}

\begin{ex} %[2D2V12]
An có 16 cái nhãn vở, Bình có ít hơn An 6 cái nhãn vở. Hỏi An phải cho Bình bao nhiêu cái nhãn vở để hai bạn có số nhãn vở bằng nhau?
\shortans{3}
\loigiai{
\textbf{Đáp số:} $3$\\
An hơn Bình 6 cái. An phải cho Bình: $6 : 2 = 3$ cái nhãn vở.
}
\end{ex}

\begin{ex} %[2D2V12]
Một giá sách tầng trên có 18 quyển sách, tầng dưới có ít hơn tầng trên 8 quyển sách. Hỏi cả hai tầng có tất cả bao nhiêu quyển sách?
\shortans{28}
\loigiai{
\textbf{Đáp số:} $28$\\
Tầng dưới: $18 - 8 = 10$ quyển. Tổng: $18 + 10 = 28$ quyển.
}
\end{ex}

\begin{ex} %[2D2V12]
Nhà Lan nuôi 15 con gà, nhà Mai nuôi nhiều hơn nhà Lan 6 con gà. Hỏi nhà Mai nuôi bao nhiêu con gà?
\shortans{21}
\loigiai{
\textbf{Đáp số:} $21$\\
$15 + 6 = 21$ con gà.
}
\end{ex}
